\documentclass[../../../include/open-logic-section]{subfiles}

\begin{document}

\olfileid[fr]{sth}{choice}{vitali}
\olsection{Annexe : le paradoxe de Vitali}

Pour juger véritablement de l'acceptabilité
de la construction de Banach--Tarski,
il faudrait examiner sa \emph{démonstration}.
Malheureusement, celle-ci exige bien plus
d'algèbre que nous ne pouvons en présenter
ici. Nous pouvons néanmoins faire quelques
remarques susceptibles de l'éclairer,\footnote{
Pour une étude beaucoup plus complète,
voir \cite{Weston2003} ou \cite{Wagon2016}.}
en nous concentrant sur le résultat suivant :

\begin{thm}[Paradoxe de Vitali (dans $\ZFC$)]
\ollabel{vitaliparadox}
Tout cercle peut être décomposé en une
infinité dénombrable de parties que l'on
peut réassembler, par rotations et
translations, pour former deux copies
de ce cercle.
\end{thm}

Le paradoxe de Vitali est beaucoup plus
facile à démontrer que celui de
Banach--Tarski. Nous l'appelons ainsi
parce qu'il découle de la construction
par Vitali, en \citeyear{Vitali1905},
d'un ensemble non mesurable. Mais les
aspects ensemblistes des deux démonstrations
sont très semblables. La différence
essentielle est que Banach--Tarski
considère une décomposition \emph{finie},
tandis que le paradoxe de Vitali en
considère une \emph{infinie dénombrable}.
Comme le dit \citet{Weston2003},
le paradoxe de Vitali est «~certainement
loin d'être aussi frappant que celui
de Banach--Tarski, mais il montre que
des paradoxes géométriques peuvent
survenir même dans des situations
simples~».

Le paradoxe de Vitali concerne une
figure bidimensionnelle : un cercle.
Nous travaillerons donc dans le plan,
$\Real^2$. Soit $\rotationsgroup$
l'ensemble des rotations, dans le sens
horaire, autour de l'origine, d'angle
$2\pi q$ radians avec
$q\in\mathbb{Q}\cap[0,1)$.\footnote{
Note de l'édition française : l'original
parle d'angles dont la mesure en radians
est rationnelle dans $[0,2\pi)$.
Ces rotations ne sont pas fermées
par composition ni par inversion :
la somme de deux angles rationnels
peut dépasser $2\pi$, qui est
irrationnel. Il faut prendre les
multiples \emph{rationnels du tour
complet} $2\pi$, modulo $2\pi$.}
Voici quelques faits algébriques
sur $\rotationsgroup$ ; la démonstration
éclairera l'énoncé si celui-ci
vous paraît obscur.

\begin{lem}\ollabel{rotationsgroupabelian}
L'ensemble $\rotationsgroup$ forme
un {groupe} abélien pour la
composition des fonctions.
\end{lem}

\begin{proof}
Notons $0_{\rotationsgroup}$ la rotation
d'angle nul. C'est l'élément neutre
de $\rotationsgroup$, puisque
$\comp{0_{\rotationsgroup}}{\rho}
= \comp{\rho}{0_{\rotationsgroup}}
= \rho$ pour tout
$\rho \in \rotationsgroup$.

Chaque élément a un inverse. Si
$\rho \in \rotationsgroup$ tourne
d'un angle $2\pi q$, alors
$\rho^{-1} \in \rotationsgroup$
tourne d'un angle
$2\pi(1-q)$, pris modulo
$2\pi$ ; ainsi
$\rho \circ \rho^{-1}
= 0_\rotationsgroup$.

La composition est associative :
$\comp{\rho}{(\comp{\sigma}{\tau})}
= \comp{(\comp{\rho}{\sigma})}{\tau}$
pour tous
$\rho, \sigma, \tau \in
\rotationsgroup$.

Elle est commutative :
$\comp{\rho}{\sigma}
= \comp{\sigma}{\rho}$
pour tous
$\rho, \sigma \in \rotationsgroup$.
\end{proof}

Nous pouvons même partager ce groupe
en deux et reconstituer le groupe
entier à partir de chacune des moitiés :

\begin{lem}\ollabel{disjointgroup}
Il existe une partition de
$\rotationsgroup$ en deux ensembles
disjoints, $\rotationsgroup_{1}$
et $\rotationsgroup_{2}$, dont chacun
permet de retrouver tout
$\rotationsgroup$ par l'application
$\rho\mapsto\comp{\rho}{\rho}$.
\end{lem}

\begin{proof}
Soit $\rotationsgroup_{1}$
l'ensemble des rotations d'angle
$2\pi q$ pour
$q\in\mathbb{Q}\cap[0,\nicefrac{1}{2})$ ;
posons
$\rotationsgroup_{2}
= \rotationsgroup\setminus
\rotationsgroup_1$.
Par un calcul élémentaire,
$\Setabs{\comp{\rho}{\rho}}
{\rho \in \rotationsgroup_1}
= \rotationsgroup$.
Le même résultat vaut pour
$\rotationsgroup_2$.
\end{proof}

Nous allons appliquer ce fait
à la démonstration du
\olref{vitaliparadox}.
Soit $\onesphere$ le cercle unité :
l'ensemble des points du plan à
distance exactement~$1$ de l'origine,
soit
$\Setabs{\tuple{r,s} \in \Real^2}
{\sqrt{r^2+s^2}=1}$.
Pour le partager en parties,
considérons la relation suivante
sur $\onesphere$ :
\[
	r \sim s \emph{ ssi }
	(\exists \rho \in \rotationsgroup)
	\rho(r) = s.
\]
Deux points du cercle sont donc
reliés si et seulement si l'on
peut passer de l'un à l'autre
par une rotation d'un nombre
rationnel de tours autour de
l'origine. Sans surprise :

\begin{lem}
$\sim$ est une relation d'équivalence.
\end{lem}

\begin{proof}
C'est immédiat d'après le
\olref{rotationsgroupabelian}.
\end{proof}

Invoquons maintenant l'axiome du
choix pour obtenir un ensemble
$C$ contenant exactement un
élément de chaque classe d'équivalence
de $\onesphere$ selon $\sim$.
Autrement dit, considérons une
fonction de choix $f$ sur
l'ensemble de ces classes,\footnote{
Puisque $\rotationsgroup$ est
!!{enumerable}, chaque !!{element}
de $E$ est !!{enumerable}.
Puisque $\onesphere$ est
!!{nonenumerable}, il découle
du \olref{vitalicover} et de la
\olref[card-arithmetic][simp]{kappaunionkappasize}
que $E$ est !!{nonenumerable}.
Il s'agit donc d'un usage du
choix \emph{non} dénombrable.}
\[
	E = \Setabs{\equivrep{r}{\sim}}
	{r \in \onesphere},
\]
et posons $C = \ran{f}$.
Pour chaque rotation
$\rho \in \rotationsgroup$,
l'ensemble $\funimage{\rho}{C}$
contient les points obtenus
en appliquant $\rho$ à chacun
des points de $C$. Les deux
résultats suivants montrent que
ces ensembles recouvrent tout
le cercle, sans chevauchement.

\begin{lem}\ollabel{vitalicover}
$\onesphere =
\bigcup_{\rho \in \rotationsgroup}
\funimage{\rho}{C}$.
\end{lem}

\begin{proof}
Fixons $s \in \onesphere$.
Il existe $r \in C$ tel que
$r \in \equivrep{s}{\sim}$,
c'est-à-dire $r\sim s$.
Ainsi, $\rho(r)=s$ pour
un certain
$\rho \in \rotationsgroup$.
\end{proof}

\begin{lem}\ollabel{vitalinooverlap}
Si $\rho_1\neq\rho_2$, alors
$\funimage{\rho_1}{C}
\cap\funimage{\rho_2}{C}
=\emptyset$.
\end{lem}

\begin{proof}
Supposons que
$s\in\funimage{\rho_1}{C}
\cap\funimage{\rho_2}{C}$.
Alors
$s=\rho_1(r_1)=\rho_2(r_2)$
pour certains $r_1,r_2\in C$.
Il s'ensuit que
$\rho_2^{-1}(\rho_1(r_1))=r_2$.
Comme
$\comp{\rho_1}{\rho_2^{-1}}
\in\rotationsgroup$,
on a $r_1\sim r_2$.
Puisque $C$ contient exactement
un représentant de chaque classe,
$r_1=r_2$. Par conséquent
$s=\rho_1(r_1)=\rho_2(r_1)$.
Deux rotations autour de l'origine
qui coïncident en un point du
cercle unité sont identiques ;
donc $\rho_1=\rho_2$.
\end{proof}

Appliquons maintenant au cercle
nos résultats algébriques.

\begin{lem}\ollabel{pseudobanachtarski}
Il existe une partition de
$\onesphere$ en deux ensembles
disjoints $D_1$ et $D_2$,
chacun pouvant être partagé
en une infinité dénombrable
de parties que l'on peut
faire tourner pour former
une copie de $\onesphere$.
\end{lem}

\begin{proof}
Avec $\rotationsgroup_1$ et
$\rotationsgroup_2$ du
\olref{disjointgroup}, posons :
\begin{align*}
	D_{1}
	&= \bigcup_{\rho \in \rotationsgroup_1}
	\funimage{\rho}{C} &
	D_{2}
	&= \bigcup_{\rho \in \rotationsgroup_2}
	\funimage{\rho}{C}
\end{align*}
C'est une partition de
$\onesphere$ par le \olref{vitalicover},
et $D_1,D_2$ sont disjoints
par le \olref{vitalinooverlap}.
Par construction, $D_1$
se partage en un nombre
dénombrable d'ensembles
$\funimage{\rho}{C}$,
avec
$\rho\in\rotationsgroup_1$.\footnote{
Note de l'édition française :
l'original écrit ici $\rho\in R_1$,
symbole qui n'a pas été défini ;
il s'agit de $\rotationsgroup_1$.}
On peut les faire tourner pour
former une copie de $\onesphere$,
car
$\onesphere
=\bigcup_{\rho\in\rotationsgroup}
\funimage{\rho}{C}
=\bigcup_{\rho\in\rotationsgroup_1}
\funimage{(\comp{\rho}{\rho})}{C}$
par le \olref{disjointgroup}
et le \olref{vitalicover}.
Le même raisonnement s'applique
à $D_2$.
\end{proof}
\noindent
Il en découle immédiatement
le paradoxe de Vitali :
à partir de $\onesphere$,
nous obtenons \emph{deux}
copies du cercle en le partageant
en un nombre dénombrable de
parties — les différents
$\funimage{\rho}{C}$ —
et en les déplaçant rigidement :
il suffit de faire tourner
les parties de $D_1$, puis de
translater et de faire tourner
les parties de $D_2$).
%
%
%
%
%
%
%
%
%
%
%
%
%
%

Récapitulons la stratégie.
Nous avons commencé par des
faits algébriques concernant
un groupe de rotations du plan.
Nous nous en sommes servis
pour partager $\onesphere$
en classes d'équivalence.
Nous avons ensuite obtenu un
«~paradoxe~» en utilisant
l'axiome du choix pour prendre
un représentant de chaque classe.

La démonstration de
Banach--Tarski suit exactement
la même stratégie. La différence
principale est que les faits
algébriques dont elle dépend
sont beaucoup plus compliqués
que ceux de notre démonstration
du paradoxe de Vitali.
Ces faits algébriques n'ont
toutefois rien à voir avec
l'axiome du choix. Résumons-les
rapidement.

Pour démontrer Banach--Tarski,
on établit d'abord un analogue
du \olref{disjointgroup} :
tout \emph{groupe libre à au moins
deux générateurs} peut\footnote{
Note de l'édition française :
l'original dit «~tout groupe libre~».
Le groupe libre sur un générateur
est isomorphe à $\mathbb{Z}$ et
n'admet pas cette décomposition
paradoxale ; il faut au moins
deux générateurs.}
être partagé en quatre parties
que l'on peut, intuitivement,
«~déplacer~» de manière à
reconstituer deux copies
du groupe entier.\footnote{
Le fait que \emph{quatre}
parties suffisent est dû à
\cite{Robinson1947}.
Pour une démonstration récente,
voir \citet[Theorem 5.2]{Wagon2016}.
Nous suivons \citet[p.~3]{Weston2003}
en parlant du «~déplacement~»
des parties du groupe.}
On montre ensuite que deux
rotations particulières autour
de l'origine de $\Real^3$
engendrent un groupe libre
de rotations $F$.\footnote{
Voir \citet[Theorem 2.1]{Wagon2016}.}
(Nous n'avons toujours pas
utilisé le choix.)
Deux points de la sphère
sont alors dits «~semblables~»
si une rotation de $F$
fait passer de l'un à l'autre.
Nous \emph{utilisons l'axiome
du choix} pour choisir exactement
un point dans chaque classe
d'équivalence de points
«~semblables~».
En appliquant notre partage de $F$, comme
dans le \olref{pseudobanachtarski}, nous obtenons une décomposition
paradoxale de la sphère privée des points
fixes des rotations non triviales de $F$.
Ces points sont dénombrables ; un argument d'absorption permet
de les réintégrer et d'obtenir deux copies
de la sphère entière.\footnote{Note de l'édition française : l'original transfère directement les quatre parties du groupe à toute la sphère. L'action n'est libre qu'après retrait des points fixes ; leur réintégration exige un argument supplémentaire et ne conserve pas automatiquement ce nombre de parties.}
On obtient ainsi
\citep{Hausdorff1914} :

\begin{thm}[Paradoxe de Hausdorff
(dans $\ZFC$)]
La surface de toute sphère
peut être décomposée en un
nombre fini de parties que
l'on peut réassembler,
par rotations et translations,
pour former deux copies
disjointes de cette sphère.
\end{thm}

Quelques astuces algébriques
supplémentaires sont nécessaires
pour obtenir le théorème
complet de Banach--Tarski,
qui porte aussi sur l'intérieur
de la sphère. Mais, franchement,
ce ne sont que des finitions
algébriques. Weston écrit :
\begin{quote}
	[\ldots] le résultat sur
	les groupes libres constitue
	l'\emph{étape décisive}
	de la démonstration du paradoxe
	de Banach--Tarski. De ce
	point de vue, ce paradoxe
	porte moins sur $\Real^3$
	que sur la complexité
	du groupe [des translations
	et des rotations de $\Real^3$].
	\cite[p.~16]{Weston2003}
\end{quote}
Autrement dit, le fait de pouvoir
obtenir une décomposition
\emph{finie}, comme dans
Banach--Tarski, ou seulement
\emph{infinie dénombrable},
comme dans le paradoxe de Vitali,
tient à certaines propriétés
des groupes utilisés en
dimension deux ou trois.

Reconnaissons que cette dernière
observation gâche un peu la
plaisanterie à la fin de la
\olref[banach]{sec}.
Puisqu'elle est bidimensionnelle,
l'inscription «~Banach--Tarski~»
devrait être partagée en une
\emph{infinité} dénombrable
de morceaux pour être réarrangée
en «~Banach--Tarski
Banach--Tarski~».
Pour sauver la plaisanterie,
il faudrait écrire en trois
dimensions. Nous laissons
cela en exercice.

Une dernière remarque. Dans la
\olref[banach]{sec}, nous avons
dit que les «~parties~» de la
sphère obtenues dans la
décomposition ne pouvaient
pas être \emph{toutes mesurables},
mais ressemblaient à des
«~éparpillements infinis~»
impossibles à représenter.
Il en va de même pour les
parties obtenues grâce au
choix dans le \olref{pseudobanachtarski}.
Expliquons pourquoi.

Il nous faut encore esquisser
quelques notions préalables :
ce n'est \emph{qu'une} esquisse,
et vous pouvez consulter un
manuel à l'article
\emph{mesure}. Une mesure
sur un ensemble $X$ associe
à chaque $E$ d'une certaine
«~$\sigma$-algèbre~» de
parties de $X$ une valeur
$\mu(E)\in\Real$.
Les détails importent peu ici,
si ce n'est que $\mu$ doit
être dénombrablement additive :
la mesure d'une union
dénombrable d'ensembles
disjoints est la somme
de leurs mesures, c'est-à-dire
$\mu(\bigcup_{n<\omega}X_n)
=\sum_{n<\omega}\mu(X_n)$
lorsque les $X_n$ sont
deux à deux disjoints.
Dire qu'un ensemble est
«~non mesurable~» signifie
ici qu'on ne peut le faire
entrer dans le domaine de
la mesure \emph{considérée}
tout en conservant les
propriétés requises.\footnote{
Note de l'édition française :
le terme «~mesure~» suppose
usuellement des valeurs non négatives.
Sur le cercle de mesure $1$,
les parties mesurables ont une
mesure finie entre $0$ et $1$.
La non-mesurabilité est relative
à la mesure considérée et aux
invariances en question.
%
}
Avec le groupe
$\rotationsgroup$ défini
plus haut, nous obtenons :

\begin{cor}[Vitali]
Soit $\mu$ une mesure sur
une $\sigma$-algèbre de
parties de $\onesphere$
stable sous les rotations,
telle que
$\mu(\onesphere)=1$ et
$\mu(X)=\mu(Y)$ lorsque
$X$ et $Y$ sont congruents.
Alors $\funimage{\rho}{C}$
n'est pas mesurable pour
tout $\rho\in\rotationsgroup$.\footnote{
Note de l'édition française :
nous explicitons ici la stabilité
du domaine de la mesure sous les
rotations, comprise dans l'usage
d'une mesure invariante sur le cercle.
Elle permet de passer de la
mesurabilité d'une partie à celle
de ses images par rotation.
}
\end{cor}

\begin{proof}
Raisonnons par l'absurde.
Supposons que
$\mu(\funimage{\sigma}{C})=r$
pour un certain
$\sigma\in\rotationsgroup$
et un certain $r\in\Real$.
Pour chaque
$\rho\in\rotationsgroup$,
les ensembles
$\funimage{\rho}{C}$ et
$\funimage{\sigma}{C}$
sont congruents ; donc
$\mu(\funimage{\rho}{C})=r$
pour tout
$\rho\in\rotationsgroup$.\footnote{
Note de l'édition française :
l'original quantifie ici sur
$\rho\in C$, dont les éléments
sont des points et non des
rotations. La variable doit
parcourir $\rotationsgroup$.}
D'après le \olref{vitalicover}
et le \olref{vitalinooverlap},
$\onesphere=
\bigcup_{\rho\in\rotationsgroup}
\funimage{\rho}{C}$
est une union dénombrable
d'ensembles deux à deux
disjoints. L'additivité
dénombrable impose donc
que $\mu(\onesphere)=1$
soit la somme des mesures
des $\funimage{\rho}{C}$ :
\[
	1=\mu(\onesphere)
	=\sum_{\rho\in\rotationsgroup}
	\mu(\funimage{\rho}{C})
	=\sum_{\rho\in\rotationsgroup}r.
\]
Mais si $r=0$, cette somme
vaut $0$ ; si $r>0$, elle
vaut $\infty$. Dans les
deux cas, elle ne vaut
pas $1$.
\end{proof}

\end{document}
